\documentclass[11pt]{article}
\usepackage[a4paper,margin=1in]{geometry}
\usepackage{amsmath,amssymb}
\usepackage{booktabs}
\usepackage{hyperref}
\hypersetup{hidelinks}

\title{Check of the Updated Shortcut Calculation}
\author{}
\date{6 June 2026}

\begin{document}
\maketitle

\section*{Summary}

I checked the updated TeX file against the direct finite stochastic shortcut
matrix. The leading corrected resolvent formula now has the stochastic
shortcut sign, but the correction has not been propagated consistently into
the first-passage formula. In particular, the equations labelled
\texttt{first\_pass\_prob} and \texttt{first\_pass\_prob\_2} still need the
corresponding sign update to match the corrected resolvent and the finite-matrix
benchmark.

\section*{First-passage sign check}

For a shortcut that moves probability mass
\[
  \lambda=\beta(1-q)
\]
from the self-loop \(u\to u\) to the absorbing target edge \(u\to v\), the
transient dynamics loses this mass at \(u\). Summing the corrected absorbing
propagator gives
\[
  \widetilde F_{n_0}(v,z)=
  \widetilde{\mathcal F}_{n_0}(v,z)
  +\frac{
    z\beta(1-q)\widetilde W_{n_0}(u,z)
    \left[1-\widetilde{\mathcal F}_{u}(v,z)\right]
  }{
    1+z\beta(1-q)\widetilde W_{u}(u,z)
  }.
\]
Equivalently, with
\[
  \widetilde H_{N,|u-v|}(z)
  =\frac{\beta(1-q)}{q}\,
  \frac{1}{1+z\beta(1-q)\widetilde W_u(u,z)},
\]
the prefactor in the next displayed first-passage equation should be positive:
\[
  \widetilde F_{n_0}(v,z)=
  \widetilde{\mathcal F}_{n_0}(v,z)
  +qz\widetilde W_{n_0}(u,z)
  \left[1-\widetilde{\mathcal F}_{u}(v,z)\right]
  \widetilde H_{N,|u-v|}(z).
\]

As currently displayed, the first expression has a minus denominator and the
second expression has a minus prefactor. These two forms therefore do not yet
match the corrected \(H\) definition.

\section*{Finite-matrix benchmark}

I used the direct stochastic matrix with \(N=100\), \(q=2/3\),
\(\beta=0.04\), shortcut \(u=6\to v=56\) in paper indexing, and target
\(v=56\). The check used \(n_0=1,\ldots,5\) and
\(z\in\{0.1,0.3,0.6,0.85,0.95\}\).

\begin{center}
\begin{tabular}{lc}
\toprule
Formula compared with the direct stochastic matrix & max absolute difference \\
\midrule
correct plus/plus first-passage formula & \(2.08\times 10^{-17}\) \\
updated \texttt{first\_pass\_prob} denominator sign & \(3.29\times 10^{-3}\) \\
updated \texttt{first\_pass\_prob\_2} minus prefactor & \(6.34\times 10^{-2}\) \\
\bottomrule
\end{tabular}
\end{center}

\section*{Pole comparison}

For the symmetric case \(|u-v|=N/2\), the updated expression has
\[
  \widetilde H_{N,N/2}(z)
  =
  \frac{1}{\frac{q}{\beta(1-q)}+\phi(z)}.
\]
Writing \(s=1/z\), for \(N=100\), \(q=2/3\),
\[
  \gamma_1 = 1-q+q\cos\frac{2\pi}{N}=0.9986844856,
\]
and
\[
  \alpha_1 = 1-q+q\cos\frac{\pi}{N}=0.9996710402.
\]

The largest pole \(s_1\) solves
\[
  \frac{q}{\beta(1-q)}+\phi(1/s)=0.
\]
At \(s=\gamma_1\), the \(\phi(1/s)\) term is zero, so the denominator is
\(q/[\beta(1-q)]>0\) for every finite positive \(\beta\). Thus \(s_1\) cannot
equal \(\gamma_1\) for \(0<\beta\leq 1\). Numerically,

\begin{center}
\begin{tabular}{rccc}
\toprule
\(\beta\) & \(s_1\) & \(s_1-\gamma_1\) & \(\alpha_1-s_1\) \\
\midrule
0.01 & 0.9996076314 & \(9.23\times 10^{-4}\) & \(6.34\times 10^{-5}\) \\
0.04 & 0.9994512449 & \(7.67\times 10^{-4}\) & \(2.20\times 10^{-4}\) \\
0.08 & 0.9993014753 & \(6.17\times 10^{-4}\) & \(3.70\times 10^{-4}\) \\
1.00 & 0.9987831951 & \(9.87\times 10^{-5}\) & \(8.88\times 10^{-4}\) \\
\bottomrule
\end{tabular}
\end{center}

So the proposed switch at \(s_1=\gamma_1\) does not occur in the admissible
range \(0<\beta\leq 1\) for this corrected \(H\).

\section*{Double-peak scan}

Using the exact first-passage PMF for the same \(N=100\), \(q=2/3\),
\(u=6\to v=56\), target \(v=56\) setup, the result depends on what is meant by
``double peak.'' With the report's clear visual rule
\[
  h_v/\min(h_1,h_2)\leq 0.8,
\]
with two filtered peaks and peak separation at least \(10\), the scan gives:

\begin{center}
\begin{tabular}{ccc}
\toprule
\(n_0\) & clear double peak? & largest checked \(\beta\) with clear double peak \\
\midrule
1 & yes & 0.030 \\
2 & yes & 0.030 \\
3 & yes & 0.030 \\
4 & no & -- \\
5 & no & -- \\
6 & no & -- \\
\bottomrule
\end{tabular}
\end{center}

At \(\beta_c=2q/[(1-q)N]=0.04\), no case remains clearly double-peaked under
this rule; for \(n_0=1,2,3\) the inter-peak valley is already shallow,
\[
  h_v/\min(h_1,h_2)\approx 0.887\text{--}0.891.
\]
If one only asks for two local maxima, the boundary is less strict and local
maxima persist to about \(\beta=0.060\)--\(0.062\) for \(n_0=1,2,3\), but the
shape is no longer a clear two-peak distribution.

\end{document}
